\section*{Chapitre \ref{ch:dv:convertibles-and-the-credit-equity-link} --- Obligations convertibles et lien crédit-action}

\begin{solution}{exo:dv:convertibles-and-the-credit-equity-link:1}
\omterm{def:dv:convertibles-and-the-credit-equity-link:price}{Prix de conversion} $100/2.5=40$~; \omterm{def:dv:convertibles-and-the-credit-equity-link:value}{valeur de conversion} $2.5\times36=90$~; prime $110/90-1=22.2\%$.
\end{solution}

\begin{solution}{exo:dv:convertibles-and-the-credit-equity-link:2}
Le porteur peut toujours conserver l'obligation et recevoir ses coupons et son principal (valant le \omterm{def:dv:convertibles-and-the-credit-equity-link:floor}{plancher obligataire}), ou convertir immédiatement (valant la \omterm{def:dv:convertibles-and-the-credit-equity-link:value}{valeur de conversion})~;
un droit de choisir plus tard vaut au moins chacune des deux options, de sorte que l'\omterm{def:dv:convertibles-and-the-credit-equity-link:cb}{obligation convertible} vaut au moins leur maximum.
\end{solution}

\begin{solution}{exo:dv:convertibles-and-the-credit-equity-link:3}
Lorsque la \omterm{def:dv:convertibles-and-the-credit-equity-link:value}{valeur de conversion} dépasse le prix de remboursement anticipé, l'émetteur peut forcer la conversion et supprimer ainsi la valeur temps résiduelle du porteur~; celui-ci est
vendeur de cette option d'achat. L'effet est le plus marqué là où la valeur temps reste importante et où le remboursement peut être exercé~: au voisinage et au-dessus du déclencheur, ici 4.7 pour
52 contre 3.9 pour 40.
\end{solution}

\begin{solution}{exo:dv:convertibles-and-the-credit-equity-link:4}
Sous la mesure de valorisation, le titre doit rapporter en moyenne $r-q$. Avec une intensité $\lambda(S)$, il chute à zéro, ce qui représente une perte espérée de $\lambda(S)$ par
unité de temps~; la dérive de la partie diffusive doit la compenser~: $\E[dS/S]=(r-q+\lambda)\,dt-\lambda\,dt=(r-q)\,dt$.
\end{solution}

\begin{solution}{exo:dv:convertibles-and-the-credit-equity-link:5}
Un spread de crédit plus large abaisse le \omterm{def:dv:convertibles-and-the-credit-equity-link:floor}{plancher obligataire}, mais une partie de la valeur de l'\omterm{def:dv:convertibles-and-the-credit-equity-link:cb}{obligation convertible} provient du droit de conversion, qui vaut davantage que la perte
du plancher dans de nombreux scénarios (et profite de la dérive plus élevée de l'action dans le modèle)~; seule la partie de type obligataire porte la sensibilité pleine~: 0.0167 par point de base
contre 0.0296.
\end{solution}

\begin{solution}{exo:dv:convertibles-and-the-credit-equity-link:6}
Dans le \omterm{def:dv:convertibles-and-the-credit-equity-link:e2c}{modèle action-crédit}, une baisse de l'action relève l'intensité de défaut et abaisse le \omterm{def:dv:convertibles-and-the-credit-equity-link:floor}{plancher obligataire}, de sorte que l'\omterm{def:dv:convertibles-and-the-credit-equity-link:cb}{obligation convertible} baisse plus vite que l'action~: le
delta intègre un canal de crédit (1.39 contre 0.84 pour 20). Couvrir avec le delta à intensité constante laisse le desk sous-couvert précisément dans les baisses qui élargissent le crédit,
ce qui le rend implicitement long de crédit.
\end{solution}

\begin{solution}{exo:dv:convertibles-and-the-credit-equity-link:7}
Avec $p=2$, le delta à 30 vaut 1.65 et la perte après couverture par delta se réduit à 1.25 pour 100 de nominal (2.55 avec $p=1.2$)~: le lien plus fort reporte davantage du mouvement de crédit sur le delta
action, que les actions vendues à découvert couvrent alors. L'\omterm{def:dv:convertibles-and-the-credit-equity-link:cb}{obligation convertible} vaut 97.32 aujourd'hui et 86.18 après le mouvement.
\end{solution}

\begin{solution}{exo:dv:convertibles-and-the-credit-equity-link:8}
Les couvertures sont locales et fondées sur un modèle~: le delta dépend du modèle du lien de crédit, le CDS couvre les petits mouvements de spread de crédit mais non le saut de défaut en totalité, et aucun des deux
ne couvre le renchérissement de l'\omterm{def:dv:convertibles-and-the-credit-equity-link:cb}{obligation convertible} par rapport à sa valeur modèle lorsque le financement se retire et que les porteurs vendent ensemble (2008).
Le portefeuille est par ailleurs vendeur de la combinaison de mouvements de l'action et du crédit en situation de détresse, où le gamma devient négatif.
\end{solution}

\begin{solution}{pb:dv:convertibles-and-the-credit-equity-link:1}
\textbf{1.} 98.47~; \omterm{def:dv:convertibles-and-the-credit-equity-link:floor}{plancher obligataire} 82.91~; \omterm{def:dv:convertibles-and-the-credit-equity-link:value}{valeur de conversion} 75~; prime 31.3~\%.
\textbf{2.} 1.55 action par obligation~: 1.55 million d'actions vendues à découvert pour 100 millions de nominal (un million d'obligations).
\textbf{3.} 69.07 pour un cours de l'action de 10, où l'intensité de défaut vaut 18.7~\% par an.
\textbf{4.} 4.65 par obligation~: 138.64 avec la clause de remboursement anticipé contre 143.29 sans elle.
\textbf{5.} Parce que l'intensité de défaut, et donc l'actualisation des flux de l'obligation, dépend du cours de l'action.
\textbf{6.} 0.0167 par point de base et par obligation~: 16\,700 par point de base pour la position.
\textbf{7.} 40.6 millions de protection (0.0167 par obligation sur une annuité risquée de 4.12 par unité de spread de crédit).
\textbf{8.} 2.21 pour 100 de nominal, soit 2.21 millions~: la convexité.
\textbf{9.} 3.35 millions.
\textbf{10.} 1.37 action par obligation~: 0.18 million d'actions vendues à découvert en moins.
\textbf{11.} 86.62~; une perte de 2.55 millions.
\textbf{12.} Une perte de 4.38 millions.
\textbf{13.} Un gain de 1.93 million (la protection rapporte 4.47 millions).
\textbf{14.} L'élargissement du spread de crédit coûte à lui seul 3.35, et la baisse de l'action rapporte à elle seule 2.21 de convexité~; le choix du delta vaut 1.83 (4.38 contre 2.55).
\textbf{15.} Le renchérissement de l'\omterm{def:dv:convertibles-and-the-credit-equity-link:cb}{obligation convertible} par rapport à sa valeur modèle lorsque les porteurs vendent ensemble (financement retiré), et le saut de défaut.
\textbf{16.} Le delta action-crédit, calibré sur la courbe de crédit et les options de l'émetteur~: il couvre le canal de crédit que le scénario exploite.
\textbf{17.} Au minimum la sensibilité au crédit au cours actuel, davantage si le saut de défaut au taux de recouvrement actuel est important par rapport au portefeuille~;
pas davantage que ce que le desk est prêt à payer pour le carry.
\textbf{18.} Des limites sur le scénario conjoint de baisse de l'action et d'élargissement du crédit, sur la concentration par émetteur, sur l'effet de levier et sur la part du portefeuille financée à court terme.
\textbf{19.} 2.55 pour 100 de nominal, soit 2.55 millions sur 100 millions, avec le delta action-crédit (4.38 millions avec le delta à intensité constante).
\textbf{20.} De la convexité sur l'action, achetée avec du risque de crédit et financée par de l'argent emprunté.
\end{solution}

\begin{solution}{iq:dv:convertibles-and-the-credit-equity-link:1}
À gauche~: proche du \omterm{def:dv:convertibles-and-the-credit-equity-link:floor}{plancher obligataire}, qui baisse lui-même lorsque l'action chute et que le crédit se dégrade (région de détresse). Au milieu~: convexe, à la fois obligation et option
(région hybride, où se loge l'arbitrage). À droite~: proche de la \omterm{def:dv:convertibles-and-the-credit-equity-link:value}{valeur de conversion}, de type action, plafonné près du déclencheur par la
clause de remboursement anticipé de l'émetteur.
\iqlookfor{les trois régions et le plafond imposé par le remboursement anticipé.}
\end{solution}

\begin{solution}{iq:dv:convertibles-and-the-credit-equity-link:2}
En $x=\ln S$~: $V_t+\frac12\sigma^2V_{xx}+(r-q+\lambda-\frac12\sigma^2)V_x-(r+\lambda)V+\lambda RF=0$. À l'échéance~: $\max(F+\text{coupon},\kappa S)$.
Contraintes après chaque pas~: $V\ge\kappa S$~; $V\le\max(\text{prix du remboursement},\kappa S)$ lorsque l'obligation est remboursable~; $V\ge$ prix de l'option de vente aux dates de vente~; coupons ajoutés
aux dates de coupon. Conditions aux limites~: recouvrement pour $S$ très faible, \omterm{def:dv:convertibles-and-the-credit-equity-link:value}{valeur de conversion} pour $S$ très élevé.
\iqlookfor{les termes d'intensité de défaut, et les contraintes comme projections.}
\end{solution}

\begin{solution}{iq:dv:convertibles-and-the-credit-equity-link:3}
La convexité (gamma échangé contre les actions vendues à découvert), le coupon net du coût d'emprunt de la vente à découvert, et la convergence des obligations bon marché vers leur valeur. Risques~:
crédit (spread et défaut), delta du modèle, liquidité et financement (ventes forcées), short squeeze ou rappel de titres empruntés sur les actions, et
remboursements anticipés par l'émetteur.
\iqlookfor{sources de rendement et liste des risques.}
\end{solution}

\begin{solution}{iq:dv:convertibles-and-the-credit-equity-link:4}
Faire remonter l'EDP dans le temps (Crank--Nicolson, avec des pas implicites de démarrage après chaque discontinuité), puis projeter~: prendre le maximum avec
la \omterm{def:dv:convertibles-and-the-credit-equity-link:value}{valeur de conversion}, le minimum avec le prix de remboursement (ou la conversion) lorsque la clause est active et que son déclencheur est atteint, le maximum avec le prix de l'option de vente aux
dates de vente~; ajouter les coupons à leurs dates. Traiter les déclencheurs de \omterm{def:dv:convertibles-and-the-credit-equity-link:softcall}{soft call} observés sur plusieurs jours par un état supplémentaire ou une approximation, et le documenter.
\iqlookfor{projection après chaque pas, et soin aux discontinuités.}
\end{solution}

\begin{solution}{iq:dv:convertibles-and-the-credit-equity-link:5}
Scénarios conjoints de l'action et du crédit~: actions en baisse de 10 à 40~\%, avec des spreads de crédit de 100 à 500 points de base plus larges au nouveau cours, sauts de défaut des plus grands
émetteurs, un choc de volatilité, un rappel de titres empruntés, et un scénario de liquidité dans lequel les \omterm{def:dv:convertibles-and-the-credit-equity-link:cb}{obligations convertibles} se renchérissent de plusieurs points par rapport au modèle, sans mouvement
des actions ni du crédit.
\iqlookfor{mouvements conjoints action--crédit plus liquidité.}
\end{solution}

\begin{solution}{iq:dv:convertibles-and-the-credit-equity-link:6}
Le portefeuille doit vendre environ la moitié de ses positions ou trouver un nouveau financement. Les ventes d'une même obligation par de nombreux fonds font baisser leurs cours sous les valeurs modèles~; les
couvertures (actions vendues à découvert, CDS) ne compensent pas ce renchérissement~; les pertes réduisent le capital et forcent de nouvelles ventes. Les prix ne se redressent que lorsque de nouveaux capitaux arrivent,
lentement.
\iqlookfor{ventes forcées, renchérissement par rapport au modèle et redressement lent.}
\end{solution}
